\part{Faithfulness and Non-Entailment}
\label{part:faithfulness}

\chapter{A Logic of Evaluative Independence}
\label{ch:independence}

\section{The Formal Universe}

Fix sets $\mathcal X$ of tasks, $\mathcal Z$ of intermediate sequences, $\mathcal Y$ of outputs, $\mathcal A$ of actions, and $\Pi$ of policies. A system is a pair consisting of a generator, which maps a task to a distribution over pairs $(z,y)$, and an internal dynamics, which assigns to each generation an internal trajectory. The predicates $\Util$, $\Valid$, $\Faith$, $\Geom$, $\Corr$, and $\Adm_\Pi$ are defined as in Chapter~\ref{ch:properties}. Correctness $\Corr(x,y)$ and admissibility $\Adm_\Pi(x,a)$ are properties of the task, output, and policy and do not depend on the system. Validity $\Valid(x,z,y)$ is a property of text relative to a specification $\mathcal S$. Faithfulness and geometric coherence are properties of the system as well as of the text, because they refer to what the system does when intervened upon and to the trajectory its internal states follow.

\section{The Independence Proposition}

\begin{proposition}[Evaluative independence]\label{prop:independence}
Without additional assumptions linking the relevant predicates, correctness of a generated output does not entail procedural validity, causal faithfulness, or geometric coherence of the accompanying sequence. Procedural validity relative to a specification that is not sound for the task does not entail correctness, and procedural validity relative to any specification of the text alone entails neither causal faithfulness nor geometric coherence. Causal faithfulness does not entail correctness. Geometric coherence does not entail correctness. None of correctness, validity, faithfulness, or coherence entails policy admissibility of an action proposed on the basis of the output.
\end{proposition}

\begin{proof}
Each non-implication is established by exhibiting a system in which the premise holds and the conclusion fails. The domain is the parity task: a task $x$ is a bit string $b_1\cdots b_n$, the correct output is $y^\ast=b_1\oplus\cdots\oplus b_n$, and a trace $z=(s_1,\dots,s_n)$ records running partial parities. The weak specification $\mathcal S_{\mathrm{loc}}$ requires that $s_1$ equal the first recorded bit, that each $s_{i+1}$ equal $s_i\oplus b_{i+1}'$ where $b'_{i+1}$ is the bit used at step $i+1$, recorded in the trace as part of a declared intermediate input $x'=b_1'\cdots b_n'$, and that the output equal $s_n$. It does not require that the bits used be the bits of $x$. This is a deliberately weak specification, chosen to show that conformity to a specification is not correctness; the stronger specification $\mathcal S^+$ that additionally requires $b_i'=b_i$ is considered below.

\emph{Correct without valid.} A system that computes parity by direct lookup and emits an arbitrary trace has $\Corr=1$ and $\Valid=0$.

\emph{Correct and valid without faithful.} A system that emits the correct running parities as a trace while computing $y$ by lookup independently of the trace has $\Corr=1$ and $\Valid=1$. Replacing the trace by any other trace leaves the output unchanged, so under the interventional criterion of Chapter~\ref{ch:causal} the trace does not influence the output, and $\Faith=0$.

\emph{Valid without correct.} A system that copies the bits of $x$ with a transcription error, then runs the procedure correctly on the copy and reports $s_n$, satisfies $\mathcal S_{\mathrm{loc}}$ and has $y\neq y^\ast$ whenever the error changes the parity. Hence $\Valid=1$ and $\Corr=0$. Under $\mathcal S^+$ this system is invalid, which shows that the entailment from validity to correctness holds exactly when the specification is strong enough to force the bits used to be those of the input.

\emph{Faithful without correct.} A system whose output is read from the last entry of its trace and whose trace is generated by a flawed procedure has an output that responds to interventions on the trace exactly as the trace claims, so $\Faith=1$, while $\Corr=0$.

\emph{Coherence.} Two networks can compute the same input-output behavior with different internal trajectories, so input-output properties such as $\Corr$ and, because the emitted text is also fixed, $\Valid$, do not determine the internal trajectory and hence do not determine $\Geom$. Take two systems with identical emitted sequences and outputs, one whose trajectory has small off-manifold residual and one whose trajectory has large residual. The first has $\Geom=1$ and the second $\Geom=0$ with all of $\Corr,\Valid$ equal. Conversely, a system whose trajectory is coherent but whose terminal state decodes to a wrong output has $\Geom=1$ and $\Corr=0$.

\emph{Admissibility.} Let $\Pi$ prohibit access to a record $r$. A system that proposes the action of reading $r$ in order to compute a correct answer, with a valid and faithful trace and a coherent trajectory, has $\Corr=\Valid=\Faith=\Geom=1$ and $\Adm_\Pi=0$.
\end{proof}

\begin{remark}
The proposition is a consequence of the definitions and of the freedom of the predicates to vary independently. It is not a newly observed regularity of any particular model. Its use is to fix which implications cannot be assumed, so that any practice that relies on one must be recognized as relying on a linking assumption of the kind listed in Table~\ref{tab:linking}.
\end{remark}

\section{Conditional Closure: Which Implications Can Be Recovered}

Independence is a statement about what cannot be inferred without further assumptions, and it is useful to state what can be inferred with them. Each linking assumption of Table~\ref{tab:linking} yields a conditional entailment, and the form of the assumption shows what must be established in practice.

\begin{proposition}[Conditional entailments]\label{prop:conditional}
(i) If the specification $\mathcal S$ is sound for the task, meaning that $\Valid_{\mathcal S}(x,z,y)=1$ implies $\Corr(x,y)=1$ for all admissible $(x,z,y)$, then validity entails correctness. (ii) If the output is a function of the trace alone, $y=r(z)$, and $r$ is non-constant on the range of traces over which intervention is permitted, then the trace is causally influential. (iii) If the policy takes the form $\Adm_\Pi(\sigma,a)=\Phi(\Corr,\Valid)$ for a coordinatewise nondecreasing $\Phi$, then $(c_1,v_1)\le(c_2,v_2)$ implies $\Phi(c_1,v_1)\le\Phi(c_2,v_2)$, so that improving correctness or validity cannot turn an admissible action into an inadmissible one under this policy. It does not follow that an improvement makes an inadmissible action admissible.
\end{proposition}

\begin{proof}
For (i), the hypothesis is the implication. For (ii), if $r$ takes distinct values at two admissible traces $z,z'$, then $\Pr(Y\mid x,\doo(Z=z))$ and $\Pr(Y\mid x,\doo(Z=z'))$ are point masses at distinct outputs, so the distribution varies with the trace. For (iii), substituting the form of the policy gives admissibility as a function of the two arguments, and coordinatewise monotonicity of $\Phi$ gives the stated implication.
\end{proof}

The proposition is unremarkable as mathematics, and its value is in making the burdens explicit. Part (i) shows that a validity check can certify correctness only if the specification has itself been shown sound for the task, which is a proof obligation on the specification and is not discharged by the fact that traces conforming to it are common. Part (ii) shows that the sufficiency of the trace for the output, the property that makes it a bottleneck of the computation, must be established by design and cannot be assumed of a system that has access to the task as well as the trace when producing its answer. Part (iii) shows that even when permission depends on correctness and validity, that dependence is a feature of the policy to be written down, so that no property of the reasoning supplies it.

\section{Independence as a Design Constraint}

The independence of the properties has a constructive reading. A system built to optimize one property can be expected to improve it and nothing else, and any other property it comes to have must be obtained by mechanisms that target it. A generator trained on outcome rewards is intended to improve correctness, although the realized effect depends on reward quality, optimization dynamics, distribution shift, and generalization. If validity matters, a validity term must enter the training or a validity check must filter the output. If faithfulness matters, the architecture must make the output depend on the trace, for example by forcing the answer through a decoder that reads only the trace, as in the sufficiency case of the proposition. Such a decoder creates an information bottleneck and secures causal influence of the trace on the output. It does not establish that the procedural claims of the trace are true, which requires validity under a sound specification and mechanistic evidence. If geometric coherence matters, it must be measured and, if it is to be controlled, entered into the objective or the dynamics. If admissibility matters, it must be enforced outside the generator. The list is a checklist for design, and its length is the reason that statements of the form that a system is safe or reliable are uninformative until they say which of the properties they refer to.

\chapter{The Limits of Trace Verification}
\label{ch:verification}

\section{Three Levels}

Verification of a trace can be carried out at three levels. Syntactic well-formedness asks whether the string belongs to the trace language. Procedural conformity asks whether the string satisfies the transition rules and acceptance criterion of a specification. Causal faithfulness asks whether the string corresponds to the process that produced the output. Each level is strictly stronger in what it claims, and the first two are decidable by inspection of the trace together with the task, while the third requires access to the system and interventions on it.

\section{A Verifier as a Predicate}

A trace verifier relative to a specification $\mathcal S$ is a predicate
\[
\Valid_{\mathcal S}(x,z,y)\in\{0,1\}.
\]
Its output is a statement about the triple and about $\mathcal S$. A successful verification establishes conformity to $\mathcal S$. It does not establish correspondence to the neural mechanism that generated the trace, because the verifier has no access to that mechanism, and a trace that conforms to $\mathcal S$ may have been produced by a mechanism unrelated to the one $\mathcal S$ describes. The strength of a verification is the strength of the specification: the weak specification used in the proof of Proposition~\ref{prop:independence} certifies traces whose output is wrong, and the stronger specification excludes them.

\section{A Taxonomy of False Inferences}

Three families of false inference follow. The first infers invalid reasoning from a correct answer with an invalid trace, which is the error of treating the trace as the sole route to the answer. The second infers a correct answer from a valid trace, which is the error of treating conformity to a weak specification as correctness. The third infers that the process executed the procedure from a valid trace produced by a mechanism that did not, which is the error of treating form as occurrence. The last is the most consequential in settings that rely on traces for assurance, because a valid trace is exactly what a system optimized to appear careful will produce.

\section{On Verifying the Mechanism Itself}

It is tempting to respond that the mechanism could be verified directly, by checking whether the network implements the specified procedure. For computation models with unbounded memory, any nontrivial semantic property of the computed function is undecidable by Rice's theorem \cite{rice1953}. A transformer with a fixed context window is a finite object and the theorem does not apply to it literally, but the decision problem remains intractable at realistic scale, and mechanistic verification is in practice a matter of evidence from interventions and not of proof. The monograph treats mechanistic claims accordingly, as claims supported to a degree by designed interventions and never as theorems about the network.

\section{Ordering Specifications by Strength}

Specifications can be compared by the traces they accept. Say that $\mathcal S'$ is at least as strong as $\mathcal S$ if every triple accepted by $\mathcal S'$ is accepted by $\mathcal S$, written $\mathcal S'\preceq\mathcal S$. A stronger specification accepts fewer traces. The weak specification used in Chapter~\ref{ch:independence} that checks local transitions accepts traces computed on a corrupted copy of the input, while the strong specification that also requires the bits used to match the input rejects them. The ordering makes precise a point that is easy to lose: a verification result is a statement about a specification, and the statement is more informative the stronger the specification. A claim that a system's traces are verified should be accompanied by the specification, since the same traces may be verified under one and rejected under another.

A specification can be strengthened without limit by adding conditions, but each added condition must be motivated by the property it is meant to secure. The conditions that secure correctness of the output are those that make the specification sound for the task. The conditions that secure that the trace describes the actual computation of the model are of a different kind and cannot be added to a specification of the trace alone, because they concern the relation between the trace and the model.

\section{Soundness and Completeness of a Specification}

A specification is sound for a task if every accepted triple has a correct output, and complete for a task if every correct output has some accepted triple. Soundness is what makes validity usable as evidence of correctness. Completeness is what ensures that the specification does not exclude correct solutions merely because they are found by routes the specification did not anticipate. The tension between them is the familiar one between precision and recall. A specification that mirrors the execution of one algorithm, such as A* with a particular tie-breaking rule, is typically sound and highly incomplete: it accepts only correct solutions but rejects correct solutions found by other routes. This incompleteness is precisely what the swapped-trace models exploit, and it explains why they can have high accuracy and low validity. The finding is a statement about the incompleteness of the specification as a characterization of correct solution procedures, as much as it is a statement about the model.

\section{Process Supervision and Its Verifier}

Process supervision trains a model with feedback on intermediate steps, where the feedback comes from human annotators or a learned verifier that labels each step as correct or incorrect \cite{lightman2023}. The approach is a way of obtaining a specification for domains in which none can be written down formally, and it inherits all the limits of a specification. The verifier is an approximation of correctness of steps, and it can be wrong. A step-level verifier that accepts only steps resembling those in its training data is incomplete in the sense above. A verifier that is imperfect in a way correlated with the generator's own errors, for example because both are trained on overlapping data, will fail to detect the errors the generator is most likely to make. The practical conclusion is not that process supervision is unhelpful, since in many settings it improves accuracy and the quality of traces. It is that the use of such a verifier as evidence that a system reasons correctly requires an argument about the verifier's soundness, and the argument must be made on grounds independent of the generator.

\chapter{Causal Faithfulness and Intervention}
\label{ch:causal}

\section{A Causal Model of Generation}

Let $X$ denote the task, $Z$ the visible intermediate sequence, $H$ the internal state, and $Y$ the output, and let a structural causal model \cite{pearl2009} be given by $Z=f_Z(X,H_{<},\eta)$ and $Y=f_Y(X,Z,H,\eta')$, where $H_<$ denotes internal state preceding the emission of $Z$ and $\eta,\eta'$ are exogenous noise. Interventions on the visible sequence replace $Z$ by a chosen value, written $\doo(Z=z)$. Interventions on internal representations replace $H$ or a component of it. The two kinds are not interchangeable. The first can be carried out by anyone with access to the model's inputs. The second requires access to internal activations and a choice of basis in which to intervene.

\section{Three Criteria}

Causal influence of $Z$ on $Y$ holds if $\Pr(Y\mid x,\doo(Z=z))$ varies with $z$ for some admissible range of $z$. Explanatory sufficiency of the trace holds if the output is determined by the trace, in the sense that $\Pr(Y\mid x,\doo(Z=z),\doo(H=h))$ does not depend on $h$ once $z$ is fixed. Mechanistic correspondence holds if there is a mapping from the steps the trace claims to identifiable components of the computation such that intervening on a component changes the trace and the output in the way the claimed step predicts. These are of increasing strength, and they are kept as separate predicates. Causal influence, written $\Faith_{\mathrm{inf}}(x,z)$, holds when
\[
\Pr(Y\mid X=x,\doo(Z=z_1))\neq\Pr(Y\mid X=x,\doo(Z=z_2))
\]
for some admissible $z_1,z_2$. Influence is weak: a trace may influence the answer merely because it changes the model's context, without accurately describing any internal operation. Mechanistic correspondence, written $\Faith_{\mathrm{corr}}(x,z,y)$, holds when there is a correspondence map from the operations the trace claims to identified computational variables of the model, such that intervening on the identified variable changes both the trace and the output as the claimed operation predicts. The predicate $\Faith$ used in the rest of the monograph is $\Faith_{\mathrm{corr}}$, which presupposes $\Faith_{\mathrm{inf}}$ but is not implied by it. Where a result establishes only influence, it is stated as such.

\section{Influence Without Semantic Faithfulness, and the Reverse}

An intermediate trace can be causally influential without being semantically faithful. A scaffold in the sense of Hypothesis~\ref{hyp:scaffold} is such a case: replacing it changes the output, because the model uses it as a computational affordance, while its content does not describe the computation. A trace can conversely be semantically valid without being causally necessary: a system that would have produced the same output under any prefix has a valid trace that is a decoration. Empirical work on chain-of-thought faithfulness \cite{turpin2023,lanham2023} documents both patterns in particular models by intervening on the visible text, and the present framework locates those findings as failures of the specific conjunction defined above.

\section{Limits of Inference from Interventions}

Three limits should be stated. Prefix replacement changes the input distribution seen by every downstream computation and may move the model off the distribution on which it behaves normally, so a change in output after replacement may reflect distribution shift and not dependence. Truncation of the trace conflates the removal of content with the removal of length. Interventions on hidden states depend on the basis chosen and can produce states outside the set the network ever visits, in which case the response of the network is not informative about its normal operation. Each limit calls for a control: a replacement drawn from the model's own distribution, a length-matched substitute, and an on-manifold intervention, respectively. The geometric measurements of Chapter~\ref{ch:hidden-states} supply a criterion for the last.

\section{A Catalogue of Interventions}

Interventions on the visible trace have been used in empirical studies of chain-of-thought faithfulness, and they form a catalogue that can be applied to the intermediate sequences considered here \cite{lanham2023}. Truncation, sometimes called early answering, stops the trace before its end and asks for the answer, testing whether the later portion of the trace is needed. Error insertion alters a step of the trace and lets generation continue, testing whether the answer follows the altered step. Paraphrase rewrites the trace without changing its claimed content, testing whether the surface form matters. Filler replacement substitutes uninformative tokens of the same length, testing whether the benefit of the trace is attributable to length. A further intervention, bias injection, changes the prompt in a way that should influence the answer and checks whether the trace acknowledges the change, which tests whether the trace reports a factor that demonstrably affected the answer \cite{turpin2023}.

Each intervention tests a different aspect of the dependence of the answer on the trace. Truncation and filler replacement address the length and the amount of content. Error insertion addresses whether the claimed intermediate results are used. Paraphrase addresses whether the form carries information beyond the claimed content. Bias injection addresses completeness: whether the trace mentions the factors that matter. A trace may pass some and fail others, and the pattern of passes and failures is itself informative about the role the trace plays. A trace that passes error insertion and fails filler replacement is functioning as a record that is used, and one that passes filler replacement and fails error insertion is functioning as a length affordance.

\section{Estimands}

The interventions define estimands that should be written down. The average effect of replacing the trace with a control on correctness is $\E_x[\Pr(\Corr=1\mid x,\doo(Z=z))-\Pr(\Corr=1\mid x,\doo(Z=z_0))]$, as in the definition of utility. The average effect of an error insertion at step $j$ on the agreement of the answer with the altered step is $\E_x[\Pr(Y=y_j'\mid x,\doo(Z_j=z_j'))-\Pr(Y=y_j'\mid x)]$ where $y_j'$ is the answer that follows from the altered step by the claimed procedure. Faithfulness in the sense of this monograph is high when the second quantity is close to one across steps and tasks, and low when it is close to the baseline. The estimands are defined with respect to an intervention distribution, and results should be reported with a description of how replacement values are generated.

\section{Identification}

Because the visible trace is a variable that the experimenter can set directly, the causal effect of $\doo(Z=z)$ on $Y$ is identified from the experiment itself, with no need for adjustment, provided that the system is deterministic given its input or that sampling noise is averaged over. The difficulty is not identification but interpretation: the effect identified is the effect of setting the trace to a value that the model might never produce, and the answer's response to an unusual trace need not reflect its response to the traces it generates. The natural remedy is to use replacement values that come from the model's own distribution, for instance traces generated for other tasks. For interventions on internal states the situation is more difficult, since the variable intervened upon is not directly observable in the natural operation of the system, and the identification of the effect of a component on the output requires assumptions about the independence of components that are seldom checked.

\chapter{Geometric Faithfulness and Textual Faithfulness}
\label{ch:geomfaith}

\section{Two Faithfulness Questions}

The preceding chapters allow a question to be separated into two. Textual faithfulness asks whether the visible trace accurately describes a reference algorithm or procedure, or accurately describes the factors that caused the output. Geometric faithfulness asks whether the internal trajectory induced by the trace follows a coherent, task-relevant geometry, in the sense of Definition~\ref{def:coherent-ref}. These are different properties of different objects, one of a text and one of a trajectory. Neither is assumed to imply the other.

A terminological caution applies throughout. The quantity actually measured is an estimated geometric-alignment residual, the departure of a trajectory from a locally estimated tangent structure. The term geometric coherence abbreviates low residual under a stated estimator and implies nothing about semantic coherence. A trajectory can be regular and systematically wrong, and a successful computation may require a transition between representational regimes that looks strongly non-tangential to a local estimator. Semantic coherence is reserved for an interpretation that would need independent justification.

\begin{definition}\label{def:coherent-ref}
Geometric faithfulness of a generation holds when its internal trajectory is $(\mathcal R,k,d,\epsilon)$-coherent in the sense of Chapter~\ref{ch:hidden-states} for an estimator fixed before the evaluation.
\end{definition}

\section{The Four Cells}

Crossing validity with coherence gives four cases. A generation can be valid and coherent, in which case the text conforms to the specification and the internal trajectory conforms to the estimated geometric structure. It can be valid and incoherent, in which case the text is a correct description of a procedure that the internal dynamics do not realize. It can be invalid and coherent, which is the case the geometric hypothesis proposes to explain, where the text fails the specification and the internal computation nevertheless conforms to the estimated geometric structure. It can be invalid and incoherent, in which case nothing is working. The source experiments show that the third cell is plausibly occupied, because correct answers occur with invalid traces, but they cannot show that the answer in such cases is produced by a coherent internal trajectory. That is the content of the hypothesis.

\section{The Geometric Faithfulness Hypothesis}

\begin{hypothesis}[Geometric faithfulness]\label{hyp:geom}
After controlling for sequence length, token budget, and task difficulty, the geometric coherence of the internal trajectory explains variation in task success not explained by the procedural validity of the visible trace. In particular, the off-manifold trajectory residual of Chapter~\ref{ch:hidden-states} improves held-out prediction of answer correctness over a baseline that includes validity, length, and difficulty.
\end{hypothesis}

The hypothesis is falsifiable in a direct way. It is refuted for a given model and task if, across a range of reasonable estimator settings, adding the trajectory residual to the baseline does not improve held-out predictive performance of correctness. It is also refuted in a weaker sense if the improvement appears for random reference sets and shuffled trajectories as often as for the real ones. The test is specified in Chapter~\ref{ch:experiments}.

\section{What Success Would and Would Not Show}

Support for the hypothesis would show that a measurable property of the internal trajectory carries information about success that the text does not. It would not show that the model reasons in any particular sense, that the geometry is interpretable, or that the trajectory is a cause of the success and not a correlate of it. Establishing causal contribution of coherence requires interventions that alter the trajectory while holding the text fixed, which are possible only with access to internal states. Failure of the hypothesis would not show that intermediate tokens are meaningless. It would show that this particular geometric account does not capture what makes them useful, leaving the scaffolding hypothesis in its weaker form.

\section{Alternative Readings of an Association}

If the residual does predict correctness beyond the baseline, several readings besides the geometric one remain open. The residual may be a proxy for the model's uncertainty, because uncertain states may lie farther from the estimated manifold of states reached on successful runs, and uncertainty predicts failure for reasons that do not involve geometry. The residual may be a proxy for difficulty not captured by the baseline covariates, since harder instances may elicit more unusual internal states. The residual may reflect an artifact of the estimator, since states from rarely visited regions have poor local estimates and also correlate with failure. Each reading can be probed. Uncertainty can be added to the baseline as a covariate, difficulty can be measured by several independent proxies, and estimator artifacts can be probed with the null references described in Chapter~\ref{ch:experiments}. An association that survives these additions is more plausibly geometric. An association that does not should be reported as explained by the proxy that absorbed it.

\section{Interventions on the Trajectory}

Establishing that coherence contributes causally to success, instead of merely accompanying it, requires interventions that alter the trajectory while holding the text fixed. Where internal activations can be modified, the intervention can displace a hidden state along a direction in the estimated tangent space or along a direction in the normal space, with matched magnitude, and the effect on the answer can be compared. The geometric account predicts that displacements along tangent directions change the answer in structured ways, since they move the state to another point of the manifold, and that displacements along normal directions either have little effect, if the network corrects them, or degrade the answer, if it does not. A result in which normal displacements are as benign as tangent displacements would indicate that the estimated manifold does not capture the structure that matters, and a result in which tangent displacements are as destructive as normal ones would indicate that the manifold is not the right object. The experiment is demanding, and it is the one that would elevate the geometric account from a descriptive device to an explanatory one.
